\section{Primitive inputs to the record and report bounds}
\label{app:primitive-inputs}

The record argument must reproduce a commitment to private bytes.
The report argument must reproduce a tag for one immutable HMAC row.
We give the distinguishing bounds and parameter conventions for each
comparison here. The implementation uses SHA3 and HMAC-SHA256. The
games specify which public hash and compression oracles are idealized.

\subsection{Hidden-input random-oracle replacement}

For a hidden-input entropy parameter $n$, define
\[
 \Phi_n(d,Q)=\min\!\left\{1,
 4\sqrt{2d^2Q\,2^{-n}}+4\sqrt{2Q^2d\,2^{-n}}\right\}.
\]
Don, Fehr and Huang's Theorem~2~\cite{dfh2022} supplies this bound when
private oracle inputs have point probability at most $2^{-n}$ and distinct
effective private queries encode distinct inputs for every secret value.
Exact repeats are cached before that interface is queried. The two budgets
count private classical queries and all ordinary public-oracle queries;
public queries may be quantum.

\paragraph{Per-nonce confinement.}
Use the typed image/intent interface of
Section~\ref{sec:opaque-rendering}, after its structural-equality comparison
and with its declared acquisition stop. For fresh seal $i$, the image
payload, local encryption and intent-authentication keys, and non-nonce
header fields are fixed before the independent draw
$N_i\gets\{0,1\}^{192}$. The surrounding computation accesses this image
family through its immutable handles, available metadata, image head and
five intent-digest interfaces. Copies, exact retries and pending-target
application retain the same family. Other records and downstream encodings
use the returned commitments and their own available inputs; their raw
inputs are computable without $N_i$, $B_i$ or $W_{i,P}$. This is the
per-nonce confinement condition for the comparison. It permits adaptive
creation and prefixes that depend on earlier replies, with every later
check using the replies of its current execution.

For the native digest framing, put
\[
 \operatorname{fr}_d(X)=
 \operatorname{BE64}(|d|)\Vert d\Vert
 \operatorname{BE64}(|X|)\Vert X.
\]
With the fixed recipe and local keys supplied as available encoding
parameters, the private maps for image $i$, for
$\nu\in\{0,1\}^{192}$, are
\[
 \begin{split}
  f_i(\nu,\mathsf{head})&=\operatorname{fr}_{d_0}(B_i(\nu)),\\
  f_i(\nu,(d,P))&=\operatorname{fr}_{d}(W_{i,P}(\nu)),
 \end{split}
\]
where $d_0$ is the image-head domain and $d$ ranges over the five
intent-digest domains. The real interface returns
$H(f_i(N_i,\cdot))$. The complete prefix $P$ includes the image-head
reply actually returned in that execution. For any fixed effective query,
its encoded input contains $\nu$ verbatim at a known offset, giving point
probability at most $2^{-192}$. Distinct domains have disjoint framed
inputs. Within one intent domain, different tags differ at the start of
$P$; equal tags fix its length, so distinct prefixes differ before the
target image for every $\nu$. These facts separate distinct effective
queries for every nonce value. Exact repeats retain the original cache entry.

For image $i$, let $K_i^{\rm hdr}$ count original intent prefixes and let
$d_i\le1+5K_i^{\rm hdr}$ count its distinct private digest inputs, with
$d_i=0$ for known, atom-free rows. Order the hybrids by fresh sealing
invocation. In hop $i$, earlier designated families use their ideal cached
replies and later families use locally sampled independent nonces and
their real hashes. The reduction runs the complete prefix, supplies the
fixed encoding parameters when family $i$ is created, and routes only that
family's private digest queries to its challenge interface. Confinement
lets it compute the surrounding execution from available inputs and the
actual returned replies. Local image encryption and intent authentication
introduce no additional shared SHA3 calls. The cost for family $i$ is
$\Phi_{192}(d_i,Q_i)$, where $Q_i$ counts the entire derived distinguisher,
including other real rows, downstream computations and the final test.
A never-created slot contributes zero. Summing gives
\[
 \epsilon_{\rm com}
 \le\min\left\{1,\sum_{i=1}^{M}\Phi_{192}(d_i,Q_i)\right\}.
\]
The hybrids retain the independent nonce distribution, including numerical
collisions; other real hash queries that reach a selected hidden input are
included in $Q_i$. The structural-equality comparison has its separate
collision allowance in Corollary~\ref{cor:real-to-opaque}.

The native paths mapped in Section~\ref{sec:native-correspondence} implement
these families: \texttt{seal} encodes the bounded record payload before
sampling the nonce at image offsets $[80,104)$; \texttt{unseal\_image}
returns decoded fields and the image head; and \texttt{PendingWrite}
retains its complete original target for exact application. The mapped
fences and checkpoints encode the returned heads and intent digests with
available metadata. A sibling write reseals the retained session payload
with its own fresh nonce. These image/intent paths supply the confinement
instance used by the record comparison; their digest domains and source
locations are identified in Appendix~\ref{app:opaque-artifacts}.

\subsection{One immutable HMAC row}

For the row experiment in Section~\ref{sec:report-composition}, let
\[
 h:\{0,1\}^{512}\times\{0,1\}^{256}\longrightarrow\{0,1\}^{256}
\]
be the public compression oracle, with arguments in message-block,
chaining-state order. A 32-byte key $K$ determines the two pad blocks
\[
 I(K)=(K\Vert0^{32})\oplus\mathtt{36}^{64},\qquad
 O(K)=(K\Vert0^{32})\oplus\mathtt{5c}^{64}.
\]
Zero exponents in these encodings count bytes. The ranges are disjoint and
each encoding is injective in $K$. Cache the two evaluations at the fixed
SHA-256 initial state. All oracle-dependent preparation is executed and
counted; no oracle-correlated advice is supplied for free.

Let $L$ be the maximum message length in bytes and put
\[
 L\le2^{61}-65,\qquad
 \ell=1+\left\lceil\frac{L+9}{64}\right\rceil.
\]
The SHA-256 padding encodes $8(64+|m|)$ bits, retaining the already consumed
64-byte key-pad block. Its suffix is determined by $|m|$ and it is injective
within the stated length domain. The parameter $\ell$ bounds the remaining
row compression work after the two key-pad calls.

Write $E_{\rm NMAC}(T,q,\ell)$ for a valid NMAC distinguishing bound with two
independent 256-bit inner/outer keys, the prescribed injective padding, at
most $T$ keyed calls and $q$ public compression calls. Hosoyamada and Iwata's
Proposition~6~\cite{hmacqrom2021} gives the asymptotic form
\[
 E_{\rm NMAC}(T,q,\ell)
 =O\!\left(\sqrt{\frac{(q+T)^3\ell^5}{2^{256}}}\right)
\]
under its stated padding and query conditions. The implicit constant remains
part of that published bound. Our message-length convention retains HMAC's
consumed-block offset.

The two key-pad replacements use $\Phi_{256}(2,q+\ell T)$. The subsequent
NMAC replacement gives an admissible row bound
\[
 \epsilon_{\rm row}(T,q,L)
 =\min\!\left\{1,
 \Phi_{256}(2,q+\ell T)+E_{\rm NMAC}(T,q,\ell)\right\}.
\]
For reports, the immutable transcript restricts the ideal NMAC function
to one original input. Its first reply is therefore a uniform tag, and
exact repeats retrieve the same value.
Report nonzero-ID checks act on that reply in both games. All other real
HMAC work remains in $q$; the multi-report theorem supplies the dependency
order and its actual disclosure residual.
